\subsubsection{The Adjoint Boundary-Value Problem}\label{subsec:adjoint-boundary-value-problem}

Throughout this subsection, let $L(\bm{x},D)$ be a scalar elliptic differential operator of order $2m$ and let $B_j(\bm{x},D)$, $j=1,\ldots,m$, be boundary operators which satisfy the complementing conditions. The general theory of adjoints requires rather stringent regularity assumptions on $\Omega$ and the coefficients; for simplicity we shall assume they are of class $C^\infty$ and that $\Omega$ is bounded. We make these assumptions throughout. We shall make the additional assumption that the $B_j$ are normal. This property is defined as follows.

\begin{Definition}
\label{def:normal-boundary-operators}
The boundary operators $B_j(\bm{x},D)$ are called normal, if their orders $m_j$ are different from each other and less than or equal to $2m-1$ and if, moreover, the leading-order term in $B_j$ contains a purely normal derivative, i.e., $B_j^p(\bm{x},\bm{n})\ne0$ for every $\bm{x}\in\partial\Omega$ (here $\bm{n}$ is the unit normal to $\partial\Omega$).
\end{Definition}

The orders of the $B_j$ cover only half the values from $0$ to $2m-1$. We can add additional boundary operators $S_j$, $j=1,\ldots,m$, to fill in the missing orders. Obviously, we can do this in such a way that the extended set of boundary operators still satisfies the conditions of normality; we merely have to take $S_j$ to be the appropriate powers of $\partial/\partial n$. We make the following definition.

\begin{Definition}
\label{def:dirichlet-boundary-system}
The boundary operators $F_j(\bm{x},D)$, $j=1,\ldots,p$, are called a Dirichlet system of order $p$, if their orders $m_j$ cover all values from zero to $p-1$ and if, moreover, the leading-order term in $F_j$ contains a purely normal derivative, i.e., $F_j^p(\bm{x},\bm{n})\ne0$ for every $\bm{x}\in\partial\Omega$ (here $\bm{n}$ is the unit normal to $\partial\Omega$).
\end{Definition}

We have the following lemma.

\begin{Lemma}
\label{lem:dirichlet-system-tangential-triangular-reduction}
Let $F_i(\bm{x},D)$, $i=1,\ldots,p$, be a Dirichlet system, and suppose the order of $F_i$ is $i-1$. Then there exist tangential differential operators $\Phi_{ij}(\bm{x},D)$ and $\Psi_{ij}(\bm{x},D)$, of order $i-j$, such that
\begin{equation}\label{eq:dirichlet-system-tangential-triangular-reduction}
\begin{aligned}
F_i(\bm{x},D)&=\sum_{j=0}^i\Phi_{ij}(\bm{x},D)\frac{\partial^{j-1}}{\partial n^{j-1}},\\
\frac{\partial^{i-1}}{\partial n^{i-1}}&=\sum_{j=0}^i\Psi_{ij}(\bm{x},D)F_j(\bm{x},D).
\end{aligned}
\end{equation}

\end{Lemma}

The existence of the $\Phi_{ij}$ is obvious from the definition. The $\Psi_{ij}$ are then obtained by inverting the triangular matrix of the $\Phi_{ij}$. We leave the details of the proof as an exercise; see Problem~\ref{ex:dirichlet-system-triangular-reduction-details}.

\begin{Corollary}
\label{cor:dirichlet-system-prescribed-traces}
Let $F_i$, $i=1,\ldots,2m$, be a Dirichlet system, and let $m_i$ denote the order of $F_i$. Let $g_i\in H^{2m+k-m_i-1/2}(\partial\Omega)$ be given. Then there exists $u\in H^{2m+k}(\Omega)$ such that $F_iu=g_i$ on $\partial\Omega$.
\end{Corollary}

The proof follows immediately from the previous lemma and Theorem~\ref{thm:domain-higher-normal-trace}.
